\section{Introduction}

(was originally an unstructured abstract)



Smalltalk-76 is fifty years old this year. Described as the first of
the “modern Smalltalks” by Dan Ingalls, Smalltalk-76 abandoned the
token-stream interpreter semantics of Smalltalk-72 and -74 in favour
of a more straightforward syntax, supporting a much more efficient
implememtation. Smalltalk-76 introduced key innovations such as a
stack-based bytecode interpreter, classes as first-class objects, and
inheritance between classes. While the Smalltalk system has evolved
significantly in the fifty years since the Smalltalk-76 release, these
innovations remain as the foundations of all Smalltalk systems to
date.

Smalltalk — especially the actual Smaltlalk-76 implementation — never
really reached widespread distribu- tion: after all, the original
Smalltalk-76 system ran only on the Xerox Alto personal minicoputer,
of which only around 2000 were ever constructed. In the subsequent
decades, however, the Smalltalk-76 architecture has provided the
infrastructure for the majority of object-oriented langages in use
today. This paper will show how systems ranging from Smalltalk-76’s
lineal succecssors (Smalltalk-78, Smalltalk-80, Squeak, Pharo),
through other similar dyanmic object-oriented langauges (notably Ruby
and Python); to statically typed languages (notably Java); and
prototype-based langauges (Self, JavaScript, Lua) are all based on the
systems architecture designed fifty years ago for Smalltalk-76.

This paper is structured as follows: \ldots

